% !TeX root = ../../Book3.tex
\chapter{正规族}
\label{b3:ch:04}
\BookChapterNumber{b3:ch:04}
\begin{chapterabstract}
局部一致收敛保留全纯性和固定阶导数。Cauchy 估计把局部有界族变成等度连续族，经有限中心与对角线抽取得到 Montel 紧性；Rouché 定理又给出 Hurwitz 零点稳定性及单叶极限。球面距离统一处理极点与无穷极限，并区分两种正规性。
\end{chapterabstract}

\begin{learninggoals}
\begin{itemize}
\item\label{b3:item:04-goal-1}能区分逐点、局部一致和全域一致收敛，并核对导数极限。
\item\label{b3:item:04-goal-2}能用紧集共同上界与 Cauchy 估计证明等度连续。
\item\label{b3:item:04-goal-3}能完成两次对角线抽取，说明所得极限为同一个全纯函数。
\item\label{b3:item:04-goal-4}能判断无零性与单叶性的保留范围，并识别常数退化。
\end{itemize}
\end{learninggoals}

列 \(z^n\) 在开单位圆盘的每个内部紧集上趋于0，却在整个圆盘上保持上确界1。列 \(nz\) 的行为更剧烈：原点处固定为0，任何非零点模趋于无穷。函数列的局部尖峰与边界逃逸需要不同处理，正规族理论正是为了从可控函数列中选出仍可使用的极限。

\ProseParagraphBreak

本章使用\cref{b3:thm:02-cauchy-estimate}和\cref{b3:thm:03-rouche}。Joel L. Schiff 的《Normal Families》\cite{Schiff1993Normal}是这一主题的专著；本章先完整证明局部有界版本，共同省略点值的更强版本在\secref{b3:sec:1403}给出证明。

% !TeX root = ../../Book3.tex
\section{局部一致收敛}
\label{b3:sec:0401}

函数列逐点收敛，未必保留连续性；即使极限全纯，也未必可以逐项求导。Cauchy 积分公式给出的控制恰好说明，复分析应采用紧集上的一致收敛，而不是仅固定每个点观察。

\subsection{紧一致收敛}
\label{b3:subsec:040101}

\begin{definition}\label{b3:def:04-compact-convergence}
设 \(U\) 为开集，\(f_n,f:U\rightarrow\mathbb C\)。若对每个紧集 \(K\subset U\) 都有 \(\sup_K|f_n-f|\rightarrow0\)，则称 \(f_n\) 在 \(U\) 上\term[jubu-yizhi-shoulian]{局部一致收敛}[locally uniform convergence]于 \(f\)，也称紧一致收敛。
\end{definition}
开集可由递增紧集穷尽，例如取
\[
 K_j=\{\,z\in U\mid |z|\le j,\ \operatorname{dist}(z,\mathbb C\setminus U)\ge1/j\,\}.
\]
空的早期集合可略去，\(U=\mathbb C\) 时只取闭圆盘。每个 \(K_j\) 闭有界，且任意 \(U\) 内紧集最终包含于其中；因此在这一列紧集上检验一致收敛就足够。

\ProseParagraphBreak

函数列 \(z^n\) 在 \(\mathbb D\) 上局部一致趋于0，因为 \(|z|\le r<1\) 时误差不超过 \(r^n\)；却不在整个开圆盘上一致趋于0，其上确界始终是1。局部一致性不对趋近边界的点给共同误差。后续紧性抽取也只在这一拓扑中进行。

\ProseParagraphBreak

紧集穷尽并不是把开集改成一个闭区域。对于有孔且无界的区域，\(K_j\) 同时截断无穷远，并远离每个孔的边界；一个固定小盘最终包含其中，但随着 \(j\) 增大，仍可能出现新的边界附近部分。若某个 \(K_j\) 为空，不需要为它另设特殊函数空间。穷尽所要求的是每个内部紧集最终被包含，而不是每一步都有非空内部。

\ProseParagraphBreak

局部一致收敛也不会自动控制随指标移动的点。若 \(z_n\) 始终位于同一个紧集且 \(z_n\rightarrow z\)，则
\[
 |f_n(z_n)-f(z)|\le\sup_K|f_n-f|+|f(z_n)-f(z)|\longrightarrow0.
\]
第一项用一致收敛，第二项用极限连续。若 \(z_n\) 逃向边界，这个共同紧集不复存在。取 \(f_n(z)=z^n\)、\(z_n=1-1/n\)，便有 \(f_n(z_n)\rightarrow e^{-1}\)，而固定点处极限为0。两种量词不能交换。

\subsection{全纯极限}
\label{b3:subsec:040102}

\begin{proposition}\label{b3:prop:04-holomorphic-limit}
设 \(f_n\) 在开集 \(U\) 上全纯，且局部一致收敛于 \(f\)。则 \(f\) 在 \(U\) 上全纯。
\end{proposition}
\begin{proof}
先在每个小闭圆盘上用连续函数的一致极限得 \(f\) 连续。固定 \(\overline{D(a,r)}\subset U\)，每个 \(f_n\) 的 Cauchy 公式为
\(f_n(z)=(2\pi\mathrm{i})^{-1}\int_{|\zeta-a|=r}f_n(\zeta)/(\zeta-z)\,\mathrm{d}\zeta\)。圆周上一致收敛使积分在每个较小闭圆盘上一致趋于相同的 \(f\) 核积分。核积分可按 \(z\) 微分，因分母与零保持正距离，所以右端全纯。由 \(z\in D(a,r)\) 时的点值极限，它恰为 \(f(z)\)。小圆盘覆盖 \(U\)，因此处处全纯。
\end{proof}
这里既没有先求 \(f_n'\) 的逐点极限，也没有从实变量可微函数列的性质类推。连续性与全纯性是通过积分表示先后获得的。

\subsection{导数收敛}
\label{b3:subsec:040103}

\begin{proposition}\label{b3:prop:04-derivative-convergence}
在上一命题的条件下，对每个固定整数 \(k\ge1\)，\(f_n^{(k)}\) 在 \(U\) 上局部一致趋于 \(f^{(k)}\)。
\end{proposition}
\begin{proof}
固定紧集 \(K\subset U\)，选 \(\delta>0\)，使所有以 \(z\in K\) 为中心、半径 \(\delta\) 的闭圆盘都在 \(U\) 中。它们的并闭包含于一个紧集 \(K'\subset U\)。对全纯差 \(f_n-f\) 使用 Cauchy 估计，得到
\[
 \sup_K|f_n^{(k)}-f^{(k)}|\le k!\delta^{-k}\sup_{K'}|f_n-f|\longrightarrow0.
\]
选取同一个 \(\delta\) 是紧性的用途，因而结论对 \(K\) 上全部点一致成立。
\end{proof}
导数阶 \(k\) 必须固定；因子 \(k!\delta^{-k}\) 随阶数增长，不能由这一结论让 \(k\) 与 \(n\) 同时趋于无穷。以 \(z^n\) 为例，它在较小圆盘的一阶导数 \(nz^{n-1}\) 仍趋于0，但第 \(n\) 阶导数为 \(n!\)。

\ProseParagraphBreak

反向也有一个带归一化的积分判据。若区域中的 \(f_n'\) 局部一致趋于 \(g\)，且某个基点值 \(f_n(a)\) 收敛，则 \(f_n\) 局部一致收敛。证明时在一个紧集的有限圆盘覆盖中，用从基点到各圆盘中心的有限条固定多边形路径，再在盘内接短线段。积分差由路径长度与导数差控制；有限路径保证同一个误差对整紧集有效。极限函数的导数为 \(g\)，由局部差商或原函数端点式即可核对。

\ProseParagraphBreak

基点值的条件不可省略，常值函数列 \(f_n=n\) 的导数全部为0，却没有有限局部极限。若定义域不连通，应每个分支固定一个值；一个基点不能控制其他分支上的加法常数。这个反向判据把“积分一次多出常数”落实为收敛中的具体条件。

\subsection{Weierstrass 定理}
\label{b3:subsec:040104}

前两项通常合称\term[weierstrass-shoulian-dingli]{Weierstrass 收敛定理}[Weierstrass convergence theorem]。
\begin{theorem}[Weierstrass]\label{b3:thm:04-weierstrass}
全纯函数列的局部一致有限极限全纯，且任意固定阶导数局部一致收敛至极限函数的相应导数。若全纯函数项级数在每个紧集上一致收敛，则其和全纯，并可任意固定阶逐项微分。
\end{theorem}
\begin{proof}
函数列结论是前两命题。函数项级数取有限部分和，它们全纯；导数部分和的局部一致极限就是和函数的导数，因此逐项微分合法。
\end{proof}
若 \(\sum_n\sup_K|g_n|<\infty\) 对每个紧集成立，Weierstrass 判别法保证所需的一致收敛。例子 \(\sum_{n\ge1}e^{-n}z^n\) 在 \(|z|<e\) 局部一致收敛；在更大区域则不能借每项整函数声称和仍整。级数收敛域是另外必须核对的数据。

紧集一致收敛还保持局部积分恒等式。若 \(f_n\rightarrow f\) 在紧集上一致，\(\gamma\) 是区域中一条固定分段光滑曲线，则
\[
 \left|\int_\gamma(f_n-f)\,\mathrm{d}z\right|
 \le L(\gamma)\sup_{\gamma}|f_n-f|\longrightarrow0.
\]
曲线的像是紧集，长度又固定，两个条件都在这里被使用。如果曲线随 \(n\) 增长或靠近边界，上式右端未必趋于0。对单位圆盘中的 \(z^n\)，紧集一致趋于0，但在 \(|z|=1\) 上模恒为1，不能把内点极限自动延伸到闭圆盘。

\ProseParagraphBreak
另一个有用例子是 \(f_n(z)=n\sin(z/n)\)。Taylor 展开在每个固定圆盘上给 \(f_n(z)=z+O(1/n^2)\)，误差常数依赖于圆盘半径。因此 \(f_n\) 紧集一致趋于恒等映射，所有阶导数也有相应极限；但它们在整个平面上的增长与零点分布很不相同。这里的收敛只控制固定的观察窗口，远处的零点可以随 \(n\) 逃向无穷远。

% !TeX root = ../../Book3.tex
\section{正规族}
\label{b3:sec:0402}

研究一个函数时，积分公式恢复它的内部值；研究一族函数时，同一公式控制整个族的变化。正规性要求任意序列都有可用的子列，因而是一种紧性性质，而不是每个成员的光滑程度。

\subsection{定义}
\label{b3:subsec:040201}

\begin{definition}\label{b3:def:04-normal-finite}
设 \(\mathcal F\) 是区域 \(U\) 上的全纯函数族。若任取 \(f_n\in\mathcal F\)，都能抽取在 \(U\) 上局部一致收敛至有限全纯函数的子列，则称 \(\mathcal F\) 在有限值意义下是\term[zhenggui-zu]{正规族}[normal family]。
\end{definition}
极限不必属于原族，所以这里对应相对序列紧性。族 \(\{z^n:n\ge1\}\) 在圆盘上正规，其某些子列极限为0，但0原本不在该族内。族 \(\{nz:n\ge1\}\) 在包含0的圆盘上不正规：若局部一致有有限极限，导数在0也应收敛，而相应导数为无界整数。

\ProseParagraphBreak

有些教材允许子列趋于无穷。为了处理极点，下面改用球面距离统一表述；有限值版本和球面版本不能未经说明互换。

\subsection{序列紧性}
\label{b3:subsec:040202}

局部一致拓扑可以由穷尽紧集定义度量
\[
 \mathrm{d}_{\mathrm{loc}}(f,g)=\sum_{j=1}^{\infty}2^{-j}\min\{1,\sup_{K_j}|f-g|\}.
\]
级数始终收敛；三角不等式从每个截断上确界得到。该度量趋于0当且仅当所有 \(K_j\) 上一致趋于0：前向由第 \(j\) 项的下界得到，反向先使尾和小，再控制有限个前项。因此“抽子列”与这个度量空间的序列紧性相对应。

\ProseParagraphBreak

对角线抽取的顺序不能省略。先在 \(K_1\) 选收敛子列，再在其内选 \(K_2\) 上的子列，依次进行，最后取第 \(j\) 列的第 \(j\) 个成员。对任何固定 \(K_m\)，对角线尾部都属于第 \(m\) 次子列，所以仍一致收敛；各紧集上的极限在重叠处一致，拼成 \(U\) 上同一个函数。若只对每点分别抽子列，就不能保证得到了同一条子列。

\ProseParagraphBreak

局部正规性能够拼成全域正规性。
\begin{proposition}\label{b3:prop:04-normal-locality}
设 \(\mathcal F\) 为开集 \(U\) 上的全纯函数族。若每点都有邻域，使限制族在该邻域有限值正规，则整个族在 \(U\) 上有限值正规。
\end{proposition}
\begin{proof}
平面第二可数性使这些邻域有可数子覆盖 \(V_1,V_2,\ldots\)。对任意函数序列，先在 \(V_1\) 抽取局部一致子列，再依次在 \(V_j\) 抽取，最后用对角线得到每个 \(V_j\) 上局部一致收敛的同一子列。重叠处的点值极限唯一，所以局部极限拼成 \(U\) 上全纯函数。若 \(K\subset U\) 紧，可用有限个位于各 \(V_j\) 内的小闭盘覆盖 \(K\)，这些盘上的一致收敛给出全 \(K\) 一致收敛。
\end{proof}
可数性与紧性承担不同作用：可数子覆盖允许一次对角线抽取，紧集的有限子覆盖把局部误差合成一个全紧集误差。只说“每点可以抽子列”而不经过这两个步骤，还没有得到正规性。

\subsection{球面度量}
\label{b3:subsec:040203}

扩充平面的\term[qiumian-juli]{球面距离}[chordal distance]取为
\[
 \chi(z,w)=\frac{|z-w|}{\sqrt{1+|z|^2}\sqrt{1+|w|^2}},\qquad
 \chi(z,\infty)=\frac1{\sqrt{1+|z|^2}}.
\]
它是单位球面弦长的一半，经立体投影得到，故满足三角不等式。归一化常数不会改变正规性，但应固定一种写法。倒数变换满足 \(\chi(1/z,1/w)=\chi(z,w)\)，包括0和无穷远。

\ProseParagraphBreak

在有限值有界集合上，球面距离与欧氏距离等价：若 \(|z|,|w|\le M\)，则 \(|z-w|/(1+M^2)\le\chi(z,w)\le|z-w|\)。趋向无穷则表示模一致变大。函数 \(n\) 的常值列没有有限局部极限，却在球面意义下一致趋向 \(\infty\)。这一点使两种正规性明确不同。

\subsection{亚纯函数族}
\label{b3:subsec:040204}

\begin{definition}\label{b3:def:04-meromorphic-normal}
区域 \(U\) 上的亚纯函数族若每个序列都有球面局部一致收敛子列，其极限是亚纯函数或常值 \(\infty\)，则称它在球面意义下正规。
\end{definition}
\begin{proposition}\label{b3:prop:04-spherical-limit}
亚纯函数列的球面局部一致极限，若存在，则是亚纯函数或常值 \(\infty\)。
\end{proposition}
\begin{proof}
球面一致极限先是到球面的连续映射 \(f\)。若 \(f(a)\) 有限，缩小圆盘使像避开 \(\infty\) 的一个邻域；充分大的 \(f_n\) 在这个圆盘也有限且一致有界，球面收敛转成欧氏收敛，Weierstrass 定理给出 \(f\) 全纯。若 \(f(a)=\infty\)，在小圆盘中像避开0，充分大的 \(1/f_n\) 全纯有界并欧氏一致收敛到 \(1/f\)。其极限全纯。若极限不恒零，零点孤立，对倒数取回得到 \(f\) 亚纯；若它在一个圆盘恒零，\(f\) 在该盘恒为无穷。全局若同时有有限值点，则连接两种区域的边界附近，倒数极限的恒等定理迫使继续为零，产生矛盾。连通性因此使恒无穷情形传播到整个区域。
\end{proof}
族 \(nz\) 在0附近没有球面一致极限，因为点0处始终为0，任意非零点却趋于无穷，候选极限不连续。允许无穷极限仍不能消除局部尖峰。

\ProseParagraphBreak

球面一致收敛允许极点的位置移动。例子 \(f_n(z)=1/(z-1/n)\) 在每个有限平面紧集上球面一致趋于 \(1/z\)：对有限非极点，
\[
 \chi\left(\frac1{z-1/n},\frac1z\right)
 =\chi(z-1/n,z)\le\frac1n,
\]
等式利用倒数是球面等距变换，并且在移动极点处也按连续性成立。欧氏一致收敛则不可能在含0的紧集上成立，因为函数有极点。球面坐标正是为了让这种自然极限连续可见。

\ProseParagraphBreak

全纯函数列的球面极限若不是恒为无穷，则实际上仍全纯，不会凭空产生一个孤立极点。若极限在 \(a\) 为极点，取小盘使其倒数在盘内全纯、仅中心有零点。充分晚的 \(f_n\) 在小盘内远离0，所以 \(1/f_n\) 全纯无零，局部一致趋于有中心零点但不恒零的 \(1/f\)，与 Hurwitz 无零极限结论矛盾。因此极点只能来自亚纯成员或整域无穷退化。这个细节在不同正规族定义之间换用时尤其重要。

% !TeX root = ../../Book3.tex
\section{Montel 定理}
\label{b3:sec:0403}

局部有界性提供导数上界，导数上界提供等度连续性，而有限覆盖与对角线抽取产生一个共同子列。Montel 定理的有限值版本因此把解析估计转换成紧性。省略值版本需要更深的几何控制，应与这一基本版本区别。

\subsection{局部有界族}
\label{b3:subsec:040301}

每个全纯函数在内部紧集上都连续有界，但这些上界可能随函数任意增大。抽取极限所需的是控制整个族的一个共同常数。

\Needspace{6\baselineskip}
\begin{definition}\label{b3:def:04-locally-bounded-family}
设 \(U\subseteq\mathbb C\) 为开集，\(\mathcal F\) 是从 \(U\) 到 \(\mathbb C\) 的函数族。若对每个紧集 \(K\subset U\)，存在 \(M_K<\infty\)，使
\[
 |f(z)|\le M_K\quad(f\in\mathcal F,\ z\in K),
\]
则称 \(\mathcal F\) 为 \(U\) 上的\term[jubu-youjie-zu]{局部有界族}[locally bounded family]。
\end{definition}

常数可以依赖观察的紧集，不能依赖族的成员。例如 \(\{z^n:n\ge1\}\) 在单位盘上共同有界为 \(1\)，而 \(\{n:n\ge1\}\) 在任意非空集合上都没有共同界。后者每个成员仍是全纯且有界的函数，所以“每个函数有界”与“整个族共同有界”是不同的条件。

\ProseParagraphBreak

定义等价于一个点附近的共同控制：对每个 \(a\in U\)，存在邻域 \(V_a\) 及常数 \(M_a\)，使全部 \(f\in\mathcal F\) 在整个 \(V_a\) 上满足 \(|f|\le M_a\)。从紧集定义推出这一点，只要在 \(a\) 周围选一个内部闭圆盘即可；反过来，用有限个 \(V_a\) 覆盖给定紧集，把对应有限个 \(M_a\) 取最大，就得到紧集版本。这里一个邻域同时控制所有函数及邻域中的所有点，不能简化为只在中心点有共同界。

\ProseParagraphBreak

回忆等度连续的量词：对紧集 \(K\) 上的族，给定 \(\varepsilon>0\)，要选出一个 \(\delta>0\)，使任意成员 \(f\) 和任意 \(z,w\in K\) 在 \(|z-w|<\delta\) 时都满足 \(|f(z)-f(w)|<\varepsilon\)。这个 \(\delta\) 同样不能依赖成员。全纯性通过 Cauchy 导数估计把共同上界转成这样的共同连续模。

\Needspace{6\baselineskip}
\begin{lemma}\label{b3:lem:04-equicontinuity}
设 \(U\subseteq\mathbb C\) 为开集，\(\mathcal F\) 是 \(U\) 上的全纯函数族。若对每个紧集 \(L\subset U\) 有
\[
 \sup_{f\in\mathcal F}\sup_{z\in L}|f(z)|<\infty,
\]
则 \(\mathcal F\) 在每个紧集 \(K\subset U\) 上等度连续。
\end{lemma}
\begin{proof}
空紧集无须讨论。对非空紧集 \(K\subset U\)，取 \(r>0\) 足够小，使
\[
 L=\{\zeta\in\mathbb C:\operatorname{dist}(\zeta,K)\le2r\}
 \subset U.
\]
这是可能的，因为 \(K\) 与闭集 \(\mathbb C\setminus U\) 有正距离；当 \(U=\mathbb C\) 时可任取 \(r>0\)。集合 \(L\) 闭且有界，因而紧。局部有界性给出一个 \(M_L\)，控制所有成员在 \(L\) 上的模。

对任意满足 \(\operatorname{dist}(\zeta,K)\le r\) 的点，闭圆盘 \(\overline{D(\zeta,r)}\) 包含于 \(L\)。\Cref{b3:thm:02-cauchy-estimate}给出
\[
 |f'(\zeta)|\le M_L/r\quad(f\in\mathcal F).
\]
若 \(z,w\in K\) 且 \(|z-w|<r\)，连接线段上的每点与 \(K\) 的距离都小于 \(r\)，所以沿该线段积分可得
\[
 |f(z)-f(w)|
 \le \frac{M_L}{r}\cdot|z-w|.
\]
给定 \(\varepsilon>0\)，取
\(\delta=\min\{r,\varepsilon r/(M_L+1)\}\)，便有统一的等度连续估计。
\end{proof}
证明中先把 \(K\) 向外加厚，再估计导数。这样即使 \(K\) 不凸，两个邻近点的连接线段也有一个共同的可用邻域；不能仅知道 \(K\) 上的导数上界，就默认任意连接线段留在 \(K\) 中。对更高阶导数，同一方法还给出
\[
 \sup_{f\in\mathcal F}\sup_{z\in K}|f^{(k)}(z)|
 \le k!r^{-k}M_L\quad(k\ge1).
\]
以下抽取只需一阶控制，极限得到以后，每个固定阶导数的收敛由 Weierstrass 定理自动提供。

\subsection{Montel 定理}
\label{b3:subsec:040302}

\term[montel-dingli]{Montel 定理}[Montel's theorem]的基本版本把刚才的估计转成一个共同收敛子列。证明不需要先假定函数空间中的紧性定理；我们直接在紧集上完成有限网格论证。
\Needspace{6\baselineskip}
\begin{theorem}[Montel，局部有界形式]\label{b3:thm:04-montel}
设 \(U\subseteq\mathbb C\) 为区域，\(\mathcal F\) 是 \(U\) 上的局部有界全纯函数族。则任取序列 \(f_n\in\mathcal F\)，都能抽取在 \(U\) 上局部一致收敛至全纯函数的子列，即 \(\mathcal F\) 在有限值意义下正规。
\end{theorem}
\begin{proof}
任取 \(f_n\in\mathcal F\)。先固定一个非空紧集 \(K\subset U\)，证明在 \(K\) 上可抽取一致收敛子列。对每个正整数 \(m\)，紧性给出 \(K\) 的有限 \(1/m\) 网：可以选有限个 \(K\) 中的点，使相应半径 \(1/m\) 的圆盘覆盖 \(K\)。把所有网点并在一起得到一个至多可数的稠密集 \(D=\{a_1,a_2,\ldots\}\)；若 \(D\) 有限，下面只需有限次抽取。

由于 \(|f_n(a_1)|\le M_K\)，复平面中的有界数列有收敛子列。再从这个子列中选出在 \(a_2\) 处也收敛的子列，依次进行，并按上一节的办法取对角线。把所得子列记为 \(g_\nu\)。对每个固定的网点 \(a_j\)，数列 \(g_\nu(a_j)\) 都收敛。

接着把稠密点上的收敛提升为整个 \(K\) 上的一致收敛。给定 \(\varepsilon>0\)，由\cref{b3:lem:04-equicontinuity}选 \(\delta>0\)，使任意成员在距离小于 \(\delta\) 的两点之间变化小于 \(\varepsilon/3\)。稠密性和紧性允许选有限个网点 \(a_{j_1},\ldots,a_{j_s}\)，使它们的 \(\delta\) 圆盘覆盖 \(K\)。这些中心上的数列均收敛，所以存在 \(N\)，使 \(\nu,\mu\ge N\) 时
\[
 |g_\nu(a_{j_l})-g_\mu(a_{j_l})|<\varepsilon/3
 \quad(1\le l\le s).
\]
对任意 \(z\in K\)，选一个距离小于 \(\delta\) 的中心 \(a_{j_l}\)，于是
\[
 \begin{aligned}
 |g_\nu(z)-g_\mu(z)|
 &\le |g_\nu(z)-g_\nu(a_{j_l})|\\
 &\quad+|g_\nu(a_{j_l})-g_\mu(a_{j_l})|
       +|g_\mu(a_{j_l})-g_\mu(z)|\\
 &<\varepsilon.
 \end{aligned}
\]
因此 \((g_\nu)\) 在一致范数中是 Cauchy 列。对每个 \(z\)，复数域的完备性给出极限 \(g(z)\)；在最后的不等式中令 \(\mu\to\infty\)，就得到整个 \(K\) 上的一致误差界。极限在 \(K\) 上连续，且有限。

现在取 \(U\) 的递增紧集穷尽 \(K_1,K_2,\ldots\)。对原序列在 \(K_1\) 上作刚才的抽取，再对所取子列在 \(K_2\) 上抽取，依次进行。第二次对角线得到在所有 \(K_j\) 上一致收敛的同一条子列。各极限在重叠处一致，拼成 \(U\) 上的有限函数 \(f\)；每个内部紧集最终包含于某个 \(K_j\)，故收敛局部一致。最后由\cref{b3:thm:04-weierstrass}，\(f\) 全纯。
\end{proof}
两次对角线承担不同的任务：第一次在一个紧集的稠密点上取得共同子列，再用等度连续性控制其余点；第二次把一个紧集扩展到整个开集。稠密点上的逐点收敛本身没有给出函数的一致收敛，有限网点的三角估计补上的正是这一步。

\ProseParagraphBreak

例如所有全纯映射 \(f:U\to\mathbb D\) 都满足共同界 \(1\)，因此是有限值正规族。这里不必事先要求单叶；若极值问题只取其中满足单叶性及某种归一化的子族，仍然正规，因为任意成员序列也是原族中的序列。Montel 只保证极限存在并全纯，极限是否仍在所选子族中，要由下一节的零点稳定性及约束来判断。

\ProseParagraphBreak

共同点值上界也可由平均大小推出。设全纯函数族满足：对每个闭圆盘 \(\overline{D(a,R)}\subset U\)，存在 \(C_{a,R}\)，使
\[
 \int_{D(a,R)}|f(z)|^2\,\mathrm{d}A(z)\le C_{a,R}
 \quad(f\in\mathcal F),
\]
其中 \(\mathrm dA\) 为平面面积。固定这样的圆盘，对 \(z\in D(a,R/2)\) 及 \(0<t<R/2\)，圆周 Cauchy 公式和 Cauchy--Schwarz 不等式给出
\[
 |f(z)|^2
 \le\frac1{2\pi}\int_0^{2\pi}|f(z+te^{\mathrm{i}\theta})|^2\,\mathrm d\theta.
\]
乘以 \(2\pi t\)，再从 \(0\) 积到 \(R/2\)，极坐标积分给出
\[
 \pi(R/2)^2|f(z)|^2
 \le\int_{D(z,R/2)}|f(\zeta)|^2\,\mathrm dA(\zeta)
 \le C_{a,R}.
\]
所以较小圆盘上的共同界为 \(2\sqrt{C_{a,R}}/(\sqrt\pi R)\)。每点附近都有这种界，有限覆盖便给出紧集共同界，Montel 定理于是适用。这个例子说明可积的平均控制也能通过解析性转化为逐点控制。

\ProseParagraphBreak

有限值正规性也反过来迫使局部有界。
\Needspace{6\baselineskip}
\begin{proposition}\label{b3:prop:04-normal-bounded-equivalence}
设 \(U\subseteq\mathbb C\) 为区域，\(\mathcal F\) 是 \(U\) 上的全纯函数族。则 \(\mathcal F\) 在有限值意义下正规，当且仅当它在 \(U\) 上局部有界。
\end{proposition}
\begin{proof}
局部有界推出正规是 Montel 定理。反过来，假设 \(\mathcal F\) 正规却不是局部有界，则存在紧集 \(K\subset U\)，使 \(\sup_{f\in\mathcal F}\sup_K|f|=\infty\)。对每个正整数 \(n\)，可选 \(f_n\in\mathcal F\) 和 \(z_n\in K\)，满足 \(|f_n(z_n)|>n\)。

正规性给出局部一致收敛子列 \(f_{n_k}\to f\)，其中 \(f\) 全纯。因 \(f\) 在 \(K\) 上连续有界，充分大的 \(k\) 满足
\[
 |f_{n_k}(z_{n_k})|
 \le \sup_K|f_{n_k}-f|+\max_K|f|
 \le 1+\max_K|f|.
\]
这与左端大于 \(n_k\to\infty\) 矛盾。此处不必再抽取点列：整个 \(K\) 上的一致误差已经直接控制了移动点 \(z_{n_k}\)。
\end{proof}
该等价是有限值版本的结论。常值族 \(\{n:n\ge1\}\) 在球面意义正规，却不局部有界，所以不能把“有限值”三个字默默删除。教材采用哪一种定义，决定了可使用的判据强弱。

\ProseParagraphBreak

这个等价也解释了局部性。只要每个点附近都能验证族的共同上界，便可在那些邻域抽取极限，再由\cref{b3:prop:04-normal-locality}拼成全域正规性。实际使用时通常先直接估计函数；只有共同上界不易得到而子列结构已知时，才反过来使用必要性。

\subsection{省略值形式}
\label{b3:subsec:040303}

所谓共同省略一个值，指每个成员在整个定义域都不取到该值。若省略的集合还含一个固定圆盘，就能直接化为已证明的有界情形。

\Needspace{6\baselineskip}
\begin{proposition}\label{b3:prop:04-omitted-disk}
设 \(U\subseteq\mathbb C\) 为区域，\(a\in\mathbb C\)、\(\varepsilon>0\)，\(\mathcal F\) 是 \(U\) 上的全纯函数族。若
\[
 |f(z)-a|\ge\varepsilon
 \quad(f\in\mathcal F,\ z\in U),
\]
则 \(\mathcal F\) 在球面意义下正规。
\end{proposition}
\begin{proof}
任取 \(f_n\in\mathcal F\)，令 \(q_n=1/(f_n-a)\)。这些函数全纯且共同有界为 \(1/\varepsilon\)，Montel 定理给出一个局部一致收敛子列，仍记为 \(q_n\to q\)，其中 \(q\) 全纯。

定义扩充平面上的映射
\[
 T(w)=a+\frac1w,\quad T(0)=\infty,\quad T(\infty)=a.
\]
它是球面的连续映射：倒数的连续性已由等距性质得到，平移在有限点连续，在无穷远也因 \(|z+a|\to\infty\) 而连续。紧球面上的连续映射一致连续。因此普通局部一致收敛 \(q_n\to q\) 先给球面局部一致收敛，再在复合 \(T\) 后保持这种收敛，得到
\[
 f_n=T\circ q_n\longrightarrow T\circ q
\]
球面局部一致。由\cref{b3:prop:04-spherical-limit}，极限是亚纯函数或恒为 \(\infty\)，故原族球面正规；由于成员全纯，\cref{b3:cor:04-holomorphic-spherical-limit}进一步排除了非恒常的带极点极限。
\end{proof}

若 \(q\equiv0\)，上述极限就是 \(\infty\)；若 \(q\) 不恒零，它的零点原本可能在倒数中形成极点，但全纯成员的球面极限结论排除了这种可能。因而非零的 \(q\) 在整个区域无零，取回得到有限全纯极限。这一分情况说明，倒数坐标既能保留有限极限，也能记录函数整体逃向无穷远。

\ProseParagraphBreak

更强的经典结论将在\cref{b3:thm:14-montel-omitted}建立：同一区域上的全纯函数族若共同省略两个固定的不同有限复值，则球面正规；亚纯函数族若共同省略三个固定的不同球面值，也球面正规。全纯函数本来就不取 \(\infty\)，因此这两种表述对应相同的三个球面省略值。

\ProseParagraphBreak

省略一个圆盘提供了 \(|f-a|\ge\varepsilon\) 的正距离，省略有限个点只给 \(|f-a|>0\)，其下界可以随成员和自变量一起趋于零，所以刚才的倒数有界证明不够。\secref{b3:sec:1403}将通过球面导数、重缩放及三穿孔球面的覆盖结构建立强版本；这里仅说明其位置，不把它当成本章或下一章的前置。Riemann 映照定理的极值构造使用的是已经完整证明的局部有界 Montel 定理。

\ProseParagraphBreak

省略值的数量不能随意减少。整函数族 \(e^{nz}\) 始终省略 \(0\)，却在含 \(0\) 的圆盘上不球面正规：原点值固定为 \(1\)，正实点沿任意指标趋于无穷的子列都趋于 \(\infty\)，与连续极限矛盾。其在半平面上的行为留作习题。球面正规性容许常数逃向无穷，却仍要求每个内部紧集上有同一个连续极限；共同省略一个有限值本身没有排除这里的失效方式。

% !TeX root = ../../Book3.tex
\section{Hurwitz 定理与全纯极限的应用}
\label{b3:sec:0404}

局部一致收敛已经保证极限全纯，却还没有回答极限能否取得某个值、是否仍为单射。第三章的零点计数提供了进一步的控制：在一个无零边界上，极限函数的模有正下界；函数列与它的误差最终小于这个下界，Rouché 定理便把内部零点的总重数传回函数列。本节先沿这条路线研究零点和映射，再把 Montel 紧性与恒等定理结合，说明何时一组点值极限足以决定整列的局部一致极限。

\subsection{零点稳定性}
\label{b3:subsec:040401}

零点的位置可以变化，几个单零点也可以在极限中合并。稳定的量是一个固定围道内部、按重数计算的零点总数，而保证稳定的条件是极限函数在围道上不为零。\term[hurwitz-dingli]{Hurwitz 定理}[Hurwitz's theorem]把这一观察用于全纯函数列。

\Needspace{6\baselineskip}
\begin{theorem}[Hurwitz]\label{b3:thm:04-hurwitz}
设 \(U\subseteq\mathbb C\) 为区域，\(f_n:U\rightarrow\mathbb C\) 全纯，并在 \(U\) 上局部一致收敛于有限值函数 \(f\)，且 \(f\) 不恒为零。若 \(a\in U\)、\(r>0\)、\(\overline{D(a,r)}\subseteq U\)，且 \(f\) 在 \(|z-a|=r\) 上无零点，则充分大的 \(n\) 使 \(f_n\) 在该圆周无零点，并与 \(f\) 在 \(D(a,r)\) 中具有相同的零点总重数。

\ProseParagraphBreak

此外，对每个满足 \(f(a)=0\) 的点 \(a\in U\)，存在整数 \(N\ge1\) 及点 \(a_n\in U\)、\(n\ge N\)，使 \(f_n(a_n)=0\) 且 \(a_n\rightarrow a\)。
\end{theorem}
\begin{proof}
由\cref{b3:thm:04-weierstrass}，\(f\) 全纯。它不恒零，所以零点孤立，在所给闭圆盘中只有有限个零点。圆周的紧性与无零条件给出
\[
 \delta=\min_{|z-a|=r}|f(z)|>0.
\]
局部一致收敛使充分大的 \(n\) 满足
\[
 \max_{|z-a|=r}|f_n(z)-f(z)|<\delta\le |f(z)|
 \quad(|z-a|=r).
\]
于是 \(f_n\) 在圆周也无零点。对全纯主项 \(f\) 和余项 \(f_n-f\) 应用\cref{b3:thm:03-rouche}，便得圆盘内的零点总重数相同。

再设 \(a\) 是 \(f\) 的 \(m\) 阶零点。取严格递减到零的半径 \(r_j>0\)，使各闭圆盘均包含于 \(U\)，且其中 \(f\) 的唯一零点为 \(a\)。由刚证明的计数结论，对每个 \(j\) 有整数 \(M_j\)，使全部 \(n\ge M_j\) 的 \(f_n\) 在 \(D(a,r_j)\) 中均有 \(m\) 个零点，计重数。将这些最终指标增大，递归取
\[
 N_1\ge M_1,\quad
 N_{j+1}\ge\max\{M_{j+1},N_j+1\}.
\]
对每个 \(n\ge N_1\)，令
\[
 j(n)=\max\{\,j\ge1\mid N_j\le n\,\}.
\]
由于 \(N_j\) 严格增加，这个最大值存在且有限。在 \(D(a,r_{j(n)})\) 中选一个 \(f_n\) 的零点 \(a_n\)。当 \(n\rightarrow\infty\) 时，\(j(n)\rightarrow\infty\)，故 \(|a_n-a|<r_{j(n)}\rightarrow0\)。这一选择对每个足够大的指标 \(n\) 都有定义，给出了所需的整条零点列。
\end{proof}

最后这一步包含两个不同的尺度。先固定半径，才能找到适用于该圆盘的最终指标；随后让半径按层缩小，并相应推迟指标，才能把“每个小盘最终有根”写成 \(a_n\rightarrow a\)。定理不要求给同一重根分裂出的各个根作唯一编号。比如 \(z^2-1/n\) 的两个单零点 \(\pm n^{-1/2}\) 都趋于零，而极限 \(z^2\) 在零点处的重数为二。

\ProseParagraphBreak

作为一个不能直接用多项式求根公式处理的例子，考虑
\[
 f_n(z)=z^2-1+\frac{e^z}{n}.
\]
在每个固定紧集上，末项一致趋于零，所以极限为 \(z^2-1\)。对 \(a=1\) 或 \(a=-1\)，在 \(|z-a|=1/4\) 上都有
\[
 |z^2-1|=|z-a|\cdot|z+a|\ge\frac14\cdot\frac74=\frac7{16},
 \quad \frac{|e^z|}{n}\le\frac{e^{5/4}}n.
\]
因此只要 \(n>16e^{5/4}/7\)，每个圆盘 \(D(1,1/4)\)、\(D(-1,1/4)\) 中都恰有一个计重数零点，因而都是单零点。继续缩小这两个圆盘，Hurwitz 定理说明相应零点分别趋于 \(1\) 和 \(-1\)。这里的边界下界既说明为何能数根，也给出了固定圆盘适用的具体指标。

\ProseParagraphBreak

反过来，若 \(f_n(a_n)=0\)，各点 \(a_n\) 始终落在某个紧集 \(K\subseteq U\) 中，且 \(a_n\rightarrow a\in K\)，则
\[
 |f(a)|\le |f(a)-f(a_n)|+\sup_K|f-f_n|\longrightarrow0,
\]
所以聚点也是极限函数的零点。共同紧集不能省略：\(f_n(z)=1-z/n\) 在全平面局部一致趋于 \(1\)，但它的零点 \(n\) 逃向无穷远。Hurwitz 的结论只约束固定内部区域中的零点。

\ProseParagraphBreak

若极限恒零，计数便失去上述依据。函数列 \(1/n\) 无零，函数列 \(z/n\) 每个成员在零点处有一个单零点，两列却都局部一致趋于零。恒零函数没有离散零点集，也无法在任何圆周上提供严格正的最低模；不能把它看成具有某个有限总重数的函数继续套用计数结论。

\subsection{无零极限}
\label{b3:subsec:040402}

没有零点是一种在极限下几乎封闭的性质。唯一需要留下的退化情形，是整个函数同时缩小到零。

\Needspace{6\baselineskip}
\begin{corollary}[无零极限]\label{b3:cor:04-zero-free-limit}
设 \(U\subseteq\mathbb C\) 为区域，\(f_n:U\rightarrow\mathbb C\) 全纯，且对全部 \(n\ge1\)、\(z\in U\) 有 \(f_n(z)\ne0\)。若 \(f_n\) 在 \(U\) 上局部一致收敛于有限值函数 \(f\)，则 \(f\) 全纯，并且或者在 \(U\) 恒为零，或者在 \(U\) 处处无零。
\end{corollary}
\begin{proof}
全纯性由 Weierstrass 定理给出。若 \(f\) 不恒零却在 \(a\in U\) 为零，取一个闭包包含于 \(U\) 的小圆盘，使 \(a\) 是其中唯一的零点。其重数为某个正整数，边界无零。Hurwitz 定理迫使全部充分晚的 \(f_n\) 在该圆盘中也有零点，与假设矛盾。
\end{proof}

常值函数列 \(1/n\) 说明恒零例外确会发生。若 \(U\) 上的全纯函数列 \(f_n\) 共同省略某个固定值 \(b\in\mathbb C\)，且局部一致趋于有限值 \(f\)，则对 \(f_n-b\) 应用推论：\(f\) 要么恒为 \(b\)，要么仍处处不取值 \(b\)。平移到零点以后，省略值问题也变成了零点稳定性。

\ProseParagraphBreak

当极限仍无零时，倒数也有通常的局部一致极限。具体地，设 \(U\) 上的无零全纯函数 \(f_n\) 局部一致趋于处处非零的 \(f\)。固定非空紧集 \(K\subseteq U\)，有 \(c_K=\min_K|f|>0\)，且充分大的 \(n\) 使 \(|f_n|\ge c_K/2\) 于 \(K\)。于是
\[
 \sup_K\left|\frac1{f_n}-\frac1f\right|
 \le\frac2{c_K^2}\cdot\sup_K|f_n-f|\longrightarrow0.
\]
分母的共同正下界是取倒数与极限交换的依据；若 \(f_n=1/n\)，倒数则趋于无穷，不能得到有限值版本的结论。

\ProseParagraphBreak

区域的连通性同样有作用。令 \(U=D(-2,1)\cup D(2,1)\)，在左边的圆盘取 \(f_n=1/n\)，右边取 \(f_n=1\)。每个成员在 \(U\) 全纯且无零，极限却在一个分支恒零、在另一个分支无零。对一般开集，无零极限推论应逐连通分支使用。若考虑全纯函数列的球面局部一致极限，\cref{b3:cor:04-holomorphic-spherical-limit}已说明它只能是全纯函数或全域恒为无穷，不会在内点新生孤立极点。

\subsection{单叶极限}
\label{b3:subsec:040403}

区域上的单射全纯函数称为\term[danye-hanshu]{单叶函数}[univalent function]。它的不同点不能映到同一值。要研究这一性质在极限中是否保留，可以固定一个点 \(a\)，把“另一个点也映到 \(f(a)\)”改写成差商的零点问题；在中心补上导数，便得到可使用无零极限推论的全纯函数。

\Needspace{6\baselineskip}
\begin{theorem}[单叶极限]\label{b3:thm:04-univalent-limit}
设 \(U\subseteq\mathbb C\) 为区域，\(f_n:U\rightarrow\mathbb C\) 是单射全纯函数，并在 \(U\) 上局部一致收敛于有限值函数 \(f\)。则 \(f\) 是常值函数，或者是单射全纯函数。
\end{theorem}
\begin{proof}
Weierstrass 定理保证 \(f\) 全纯。假设 \(f\) 非常值，任取 \(a\in U\)，在 \(z\ne a\) 处定义
\[
 g_n(z)=\frac{f_n(z)-f_n(a)}{z-a},
 \quad g(z)=\frac{f(z)-f(a)}{z-a},
\]
并令 \(g_n(a)=f_n'(a)\)、\(g(a)=f'(a)\)。Taylor 展开说明补定义后的函数在整个 \(U\) 全纯。对 \(z\ne a\)，\(f_n\) 的单射性给出 \(g_n(z)\ne0\)；在中心处，\cref{b3:cor:01-local-injectivity}说明单射全纯函数的导数不为零，故也有 \(g_n(a)\ne0\)。

须先核对 \(g_n\rightarrow g\) 是局部一致收敛。取 \(\rho>0\)，使 \(\overline{D(a,\rho)}\subseteq U\)。在此闭圆盘上，沿直线段积分给出
\[
 g_n(z)=\int_0^1 f_n'\bigl(a+t(z-a)\bigr)\,\mathrm{d}t,
\]
\(g\) 也有对应表达。所有线段都留在同一个闭圆盘中，故
\[
 \sup_{\overline{D(a,\rho)}}|g_n-g|
 \le\sup_{\overline{D(a,\rho)}}|f_n'-f'|\longrightarrow0,
\]
其中用了\cref{b3:prop:04-derivative-convergence}。对任意紧集 \(K\subseteq U\)，其在 \(D(a,\rho/2)\) 外的部分满足 \(|z-a|\ge\rho/2\)，所以
\[
 |g_n(z)-g(z)|
 \le\frac2\rho\bigl(|f_n(z)-f(z)|+|f_n(a)-f(a)|\bigr).
\]
这部分也一致收敛，与中心附近的估计合起来即得全 \(K\) 上的一致收敛。

由于 \(f\) 非常值，\(g\) 不恒为零。对无零函数列 \(g_n\) 使用\cref{b3:cor:04-zero-free-limit}，得 \(g\) 处处无零。因而对每个 \(z\ne a\) 都有 \(f(z)\ne f(a)\)。点 \(a\) 任意，所以 \(f\) 单射。
\end{proof}

函数列 \(f_n(z)=z/n\) 在任何固定区域上都单叶，却局部一致趋于常值零。若在某个固定点 \(a\in U\) 另有 \(|f_n'(a)|\ge c>0\)，则导数收敛给出 \(|f'(a)|\ge c\)，从而排除常值退化。这里只固定 \(f_n(a)\) 的值还不够，\(z/n\) 在原点已经具有相同的函数值。

\ProseParagraphBreak

映射的收敛还提出了一个定义域问题。即使 \(f_n(z)\rightarrow f(z)\)，给定 \(w\in f(U)\) 后，方程 \(f_n(z)=w\) 是否有解仍需证明；只有先确认这一点，才谈得上 \(f_n^{-1}(w)\) 的极限。单叶全纯函数导数无零，由局部逆函数定理，其逆函数在像区域上全纯。下面证明极限像的每个紧集最终都落在这些逆函数的定义域中。

\Needspace{6\baselineskip}
\begin{proposition}[单叶逆映射的收敛]\label{b3:prop:04-inverse-convergence}
设 \(U\subseteq\mathbb C\) 为区域，\(f_n:U\rightarrow\mathbb C\) 为单叶全纯函数，并在 \(U\) 上局部一致收敛于非常值函数 \(f\)。则 \(f\) 单叶全纯，且对每个紧集 \(K\subseteq f(U)\)，存在整数 \(N\)，使全部 \(n\ge N\) 满足 \(K\subseteq f_n(U)\)，并有
\[
 \sup_{w\in K}|f_n^{-1}(w)-f^{-1}(w)|\longrightarrow0.
\]
\end{proposition}
\begin{proof}
单叶性由前一定理保证。先固定 \(w_0\in f(U)\)，记 \(a=f^{-1}(w_0)\)。取 \(r>0\)，使闭圆盘 \(C=\overline{D(a,r)}\) 包含于 \(U\)。由于 \(f\) 单射，\(f-w_0\) 在该圆盘中只有一个单零点 \(a\)，边界不为零。因此
\[
 m=\min_{|z-a|=r}|f(z)-w_0|>0.
\]
先让目标点在 \(D(w_0,m/3)\) 中移动。对其中的每个 \(w\)，常数扰动 \(|w-w_0|<m/3\) 小于圆周上的 \(|f-w_0|\)。Rouché 定理说明 \(f-w\) 在源圆盘中也恰有一个计重数零点；同时圆周上有共同下界
\[
 |f(z)-w|\ge |f(z)-w_0|-|w-w_0|>2m/3.
\]
局部一致收敛提供一个与 \(w\) 无关的最终指标，使圆周上 \(|f_n-f|<m/3\)。第二次应用 Rouché 定理，比较 \(f-w\) 与 \(f_n-w\)，得 \(f_n-w\) 在该圆盘中同样恰有一个零点。故整个 \(D(w_0,m/3)\) 同时包含于 \(f(U)\) 和这些 \(f_n(U)\)，且相应逆像都落在 \(D(a,r)\) 中。

现在在较小闭盘 \(K_0=\overline{D(w_0,m/4)}\) 上比较逆像。对 \(w\in K_0\)，记 \(z_n=f_n^{-1}(w)\)、\(z=f^{-1}(w)\)。它们都在紧集 \(C\) 内，且
\[
 |f(z_n)-f(z)|=|f(z_n)-f_n(z_n)|
 \le\sup_C|f-f_n|\longrightarrow0.
\]
连续单射 \(f:C\rightarrow f(C)\) 是从紧空间到 Hausdorff 空间的同胚，故其逆映射在紧集 \(f(C)\) 上一致连续。这里 \(f(z_n)\) 和 \(f(z)\) 均属于 \(f(C)\)，所以给定 \(\varepsilon>0\)，可选一个对全部 \(w\in K_0\) 有效的 \(\delta>0\)，使上式左端小于 \(\delta\) 时便有 \(|z_n-z|<\varepsilon\)。因而逆映射在 \(K_0\) 上一致收敛。

对给定紧集 \(K\subseteq f(U)\) 的每个点执行这一构造，所得较小像盘的内部覆盖 \(K\)。取有限子覆盖，并将各处的包含指标及一致误差指标取最大，便同时得到 \(K\subseteq f_n(U)\) 和整个 \(K\) 上的逆映射一致收敛。
\end{proof}

逆映射的定义域 \(f_n(U)\) 随指标变化，所以这里的收敛按紧集理解：先固定极限像中的一个紧集，再从某项起把所有逆映射限制到这个共同集合上。Rouché 比较中的边界正距离同时控制了目标点 \(w\) 的移动和近似函数 \(f_n\) 的误差，才使最终指标对整块像盘统一有效。

\ProseParagraphBreak

例如在单位圆盘 \(\mathbb D\) 上取
\[
 f_n(z)=\frac{z}{1+z/n}.
\]
分母在盘内无零，交叉相乘可核对 \(f_n\) 单射；在任意 \(|z|\le r<1\) 上，\(|f_n(z)-z|\le r^2/(n-r)\)，所以局部一致趋于恒等映射。解方程得到逆函数候选式
\[
 f_n^{-1}(w)=\frac{w}{1-w/n}.
\]
对 \(|w|\le r<1\)，当 \(n>r/(1-r)\) 时，这个式子的模至多为 \(r/(1-r/n)<1\)，确实给出位于原定义域中的逆像，而且
\[
 |f_n^{-1}(w)-w|\le\frac{r^2}{n(1-r/n)}\longrightarrow0.
\]
但对每个固定 \(n\)，选实数 \(n/(n+1)<w<1\)，同一逆函数候选值却大于 \(1\)，所以该 \(w\) 不属于 \(f_n(\mathbb D)\)。因此每个内部紧集最终被覆盖，并不意味着整个极限像 \(\mathbb D\) 最终都包含于 \(f_n(\mathbb D)\)。

\subsection{Vitali 型收敛判据}
\label{b3:subsec:040404}

Montel 定理提供收敛子列，还不能单独确定它们是否具有同一个极限。恒等定理恰好提供所需的唯一性：两个全纯极限只要在具有内部聚点的集合上相等，就在整个区域相等。把这两项结合，得到下面的\term[vitali-shoulian-dingli]{Vitali 收敛定理}[Vitali convergence theorem]的局部有界形式。

\Needspace{6\baselineskip}
\begin{proposition}[Vitali，局部有界形式]\label{b3:prop:04-vitali-bounded}
设 \(U\subseteq\mathbb C\) 为区域，\(f_n:U\rightarrow\mathbb C\) 全纯，并且对每个紧集 \(K\subseteq U\)，存在有限常数 \(M_K\)，使全部 \(n\ge1\) 满足 \(\sup_K|f_n|\le M_K\)。设 \(E\subseteq U\) 在 \(U\) 内有聚点，且对每个 \(z\in E\)，数列 \(f_n(z)\) 收敛于有限值。则存在全纯函数 \(f:U\rightarrow\mathbb C\)，使整列 \(f_n\) 在 \(U\) 上局部一致收敛于 \(f\)。
\end{proposition}
\begin{proof}
局部有界性使\cref{b3:thm:04-montel}适用，先从原序列抽取一条局部一致收敛子列，记其全纯极限为 \(f\)。这一步给出了候选极限的存在。

再任取原序列的一条子列。它仍属于同一个局部有界族，故 Montel 定理使它有局部一致收敛的子子列，记全纯极限为 \(g\)。对每个 \(z\in E\)，两条已收敛的子列都必须趋于原数列 \(f_n(z)\) 的点值极限，所以 \(f(z)=g(z)\)。集合 \(E\) 有位于 \(U\) 内部的聚点，区域 \(U\) 又连通，\cref{b3:thm:01-identity}因此给出 \(g=f\) 于整个 \(U\)。于是原序列的每条子列都能再抽出以 \(f\) 为局部一致极限的子子列。

最后证明整列收敛。若在某个紧集 \(K\subseteq U\) 上不一致趋于 \(f\)，则存在 \(\varepsilon>0\) 及严格增加的指标 \(n_j\)，使
\[
 \sup_K|f_{n_j}-f|\ge\varepsilon\quad(j\ge1).
\]
由前一步，这条子列又有局部一致趋于 \(f\) 的子子列，因而在 \(K\) 上的误差趋于零，与固定正下界矛盾。
\end{proof}

这一判据尤其适合从实变量极限获得复变量结论。令
\[
 F_n(z)=\left(1+\frac zn\right)^n\quad(z\in\mathbb C).
\]
每个 \(F_n\) 都是整函数，且
\[
 |F_n(z)|\le\left(1+\frac{|z|}{n}\right)^n\le e^{|z|}.
\]
任意固定紧集上的 \(|z|\) 有上界，故这给出了全族的局部共同界。在实区间 \((-1,1)\) 上，实变量的指数极限给出 \(F_n(x)\rightarrow e^x\)。Vitali 判据使整列在全平面局部一致趋于某个整函数 \(F\)；\(F\) 与 \(e^z\) 在这个实区间相等，恒等定理于是确定 \(F=e^z\)。

\ProseParagraphBreak

得到局部一致收敛后，才可以应用 Weierstrass 定理取导数。对每个固定整数 \(k\ge1\)，在 \(n\ge k\) 时有
\[
 F_n^{(k)}(z)=\frac{n(n-1)\cdots(n-k+1)}{n^k}
 \cdot\left(1+\frac zn\right)^{n-k}\longrightarrow e^z,
\]
收敛在每个紧集上一致成立。这条推导只在实区间使用了最初的点值极限，区域中的共同上界和全纯唯一性负责把它扩展到整个复平面。

\ProseParagraphBreak

局部有界条件不能删除。在 \(\mathbb D\) 上取 \(f_n(z)=(2z)^n\)，则对实区间 \(E=(-1/2,1/2)\) 中每个点都有 \(f_n(x)\rightarrow0\)，而 \(E\) 具有内部聚点；可是 \(f_n(3/4)=(3/2)^n\) 发散。该序列在包含 \(3/4\) 的紧集上没有共同上界，不能从点值数据抽出覆盖整个区域的收敛子列。

\ProseParagraphBreak

内部聚点的要求也不能改成仅在边界聚集，即使全族在区域上已有统一界。下面在单位圆盘内构造一个有界全纯函数，使它在无内部聚点的无限集合上为零。令
\[
 W(z)=\frac{1+z}{1-z},\quad
 q(z)=e^{-W(z)}-e^{-2}\quad(z\in\mathbb D).
\]
分母在 \(\mathbb D\) 无零，且
\[
 \operatorname{Re}W(z)=\frac{1-|z|^2}{|1-z|^2}>0,
\]
所以 \(q\) 全纯，并有 \(|q(z)|\le1+e^{-2}\)。对每个整数 \(k\ge1\)，取
\[
 z_k=\frac{1+2\pi\mathrm{i}k}{3+2\pi\mathrm{i}k}.
\]
其模平方为 \((1+4\pi^2k^2)/(9+4\pi^2k^2)<1\)，故 \(z_k\in\mathbb D\)；又有 \(z_k\rightarrow1\)，且各点互不相同。直接代入得 \(W(z_k)=2+2\pi\mathrm{i}k\)，从而 \(q(z_k)=0\)。于是 \(E=\{\,z_k\mid k\ge1\,\}\) 完全包含于定义域，却只在边界点 \(1\) 聚集。

\ProseParagraphBreak

现在令 \(f_{2n}=0\)、\(f_{2n-1}=q\)。这是一列在 \(\mathbb D\) 上共同有界的全纯函数；在每个 \(z\in E\) 上，整个数列 \(f_n(z)\) 恒为零，当然收敛。但 \(q(0)=e^{-1}-e^{-2}\ne0\)，所以 \(f_n(0)\) 在两个不同的值之间交替，整列连逐点收敛都不成立。本例中的两条子列极限 \(0\) 与 \(q\) 确实存在，也在 \(E\) 上相等；失效的是恒等定理所需的内部聚点，因而不能把子列的存在提升为整列收敛。


\begin{chaptersummary}
局部一致拓扑由内部紧集控制，不限制边界附近的全域上确界。Cauchy 公式使有限极限保留全纯性，并控制每个固定阶导数；Montel 定理将局部上界转换为可抽取的函数列。Hurwitz 定理保留零点总重数，继而得到无零性和单叶性，常数退化是两者都必须留下的例外。球面正规性允许无穷极限，但仍要求极限连续，不能容许内部尖峰。
\end{chaptersummary}

\BookExercises

% 本章题组自拟；以正文已建立的工具设计计算、反例、构造与综合训练。
\begin{exercise}\label{b3:ex:04-1}
证明 \(nz^n\) 在 \(\mathbb D\) 上局部一致趋于0，却不在 \(\mathbb D\) 上一致收敛。
\end{exercise}
\begin{solution}
在 \(|z|\le r<1\) 上误差不超过 \(nr^n\)，它趋于0；取足够靠近1的实数点可使 \(n|z|^n\) 任意接近 \(n\)，故全域上确界为 \(n\)。
\end{solution}

\begin{exercise}\label{b3:ex:04-2}
若 \(f_n\rightarrow f\) 在开集上局部一致且每个 \(f_n\) 全纯，证明 \(\int_\gamma f_n\,\mathrm{d}z\rightarrow\int_\gamma f\,\mathrm{d}z\)，其中 \(\gamma\) 为固定分段光滑曲线。
\end{exercise}
\begin{solution}
曲线像是内部紧集，长度估计给出两积分之差的模不超过 \(L(\gamma)\sup_\gamma|f_n-f|\)，其趋于0。曲线固定使长度不随指标增长。
\end{solution}

\begin{exercise}\label{b3:ex:04-3}
设 \(f_n\) 在圆盘全纯且 \(|f_n|\le1\)。证明存在子列，使全部阶数的原点 Taylor 系数逐个收敛，并说明为何系数极限确定同一个全纯函数。
\end{exercise}
\begin{solution}
Montel 先抽取函数局部一致收敛子列。任意固定阶导数随之收敛，所以全部阶数的系数都沿这一子列逐个收敛。Cauchy 估计给每个系数 \(|a_k|\le1\)，极限级数在盘内收敛；在较小盘上它等于函数极限的 Taylor 展开。无须对每阶再选不同子列。
\end{solution}

\begin{exercise}\label{b3:ex:04-4}
构造无零整函数列，局部一致极限恰为零；再构造在 \(\mathbb D\) 无零的函数列，其极限仍无零却趋近边界时模可任意小。
\end{exercise}
\begin{solution}
第一列取 \(1/n\)。第二列可取常列 \(1-z\)，它在开盘无零，沿 \(z=r\uparrow1\) 模趋于0。Hurwitz 的无零结论不提供统一的全域正下界。
\end{solution}

\begin{exercise}\label{b3:ex:04-5}
证明 \(\{e^{nz}:n\ge1\}\) 在右半平面球面正规，在任何包含0的圆盘不球面正规。
\end{exercise}
\begin{solution}
右半平面任意紧集上实部有正下界 \(\delta\)，故 \(|e^{nz}|\ge e^{n\delta}\) 一致趋于无穷；任意序列若指标无界可抽此类子列，指标有界则抽常子列。包含0的圆盘内，沿正负实点的候选极限分别为无穷与0，而原点固定为1，任何无界指标子列的候选极限都不连续，故不正规。
\end{solution}

\begin{exercise}\label{b3:ex:04-6}
设 \(f_n:U\rightarrow\mathbb D\) 单叶全纯，固定 \(a\in U\) 且 \(|f_n'(a)|\ge c>0\)。证明存在单叶全纯局部一致子列极限。
\end{exercise}
\begin{solution}
共同有界使 Montel 适用。导数收敛给 \(|f'(a)|\ge c\)，排除常数极限；单叶极限定理因此给出单叶性。最大模原理又使有限极限仍取值于开盘，因为若内部模达到1它为常数。
\end{solution}

\begin{exercise}\label{b3:ex:04-7}
设 \(f_n\) 全纯且局部一致趋于 \(f\)，\(f\) 在 \(a\) 有二阶零点。说明为何不能保证 \(f_n\) 在附近也有二阶零点，并给例子。
\end{exercise}
\begin{solution}
取 \(f_n(z)=z^2-1/n\)。局部一致极限为 \(z^2\)，但每个成员在 \(\pm1/\sqrt n\) 有两个单零点。Hurwitz 保留的是充分小围道内部总重数2，不是某个指定点的阶数。
\end{solution}

\begin{exercise}\label{b3:ex:04-8}
设整函数列 \(f_n\) 在每个闭圆盘共同有界，且 \(f_n(1/j)\) 对每个正整数 \(j\) 都有极限。证明整列局部一致收敛，并说明这为何不要求知道聚点0处的点值极限。
\end{exercise}
\begin{solution}
集合 \(E=\{1/j:j\ge1\}\) 在内部有聚点0。前述局部有界收敛命题直接给整列局部一致收敛。两个候选全纯极限在 \(E\) 上相同，恒等定理只需要这些零点在0聚集，不另要求预先规定0的极限。
\end{solution}
